\section{Historical Shape-Tied Control and Its Limits}
\label{app:shape_agg}
An earlier scalability experiment shared decisions among operations with
identical Timeloop coefficient rows. This is a separate restricted control;
it is not the representation used to claim coverage of the current native
master. The latter preserves independent invocation choices and deduplicates
coefficient recipes rather than silently identifying their variables.

\paragraph{Coefficient reuse versus shared decisions.}
If operations $i$ and $j$ have identical dimensions and implementations, the
same coefficient recipe may generate both sets of rows. This reduces storage
without imposing $y_i=y_j$. Replacing their selectors by one shape selector
$y_s$, however, removes solutions in which they choose different tiles.
Identical local coefficients do not imply interchangeable decisions: the
operations can have different dependencies and live-memory pressure.

\paragraph{Temporal coupling.}\label{par:shape_agg_exactness}
A product $w_{ict}=y_{ic}p_{it}$ depends on both implementation and temporal
state. A shape-indexed selector and a coarsened residency trajectory cannot
represent all per-instance products. The historical control grouped identical
Pareto rows, multiplied compute, traffic, and recomputation coefficients by
multiplicity, charged tile bytes once, and coarsened $T=N$ to $T=S$.
Decision tying restricts the domain while removing capacity rows can relax it;
there is no general ordering between its optimum and the original optimum.
An objective-preserving reduction requires an explicit lift satisfying all
original rows, not just agreement of shape dimensions.

\paragraph{Observed small-instance agreement.}
The retained comparison below solves both programs with HiGHS on the same
truncated operation set and with the original linking rows enabled. The
objectives agree to ten significant figures in these cases, with up to an
$89\times$ reduction in columns. Capacity is slack in this experiment. These
observations support the control at the tested points, not an equivalence
claim for capacity-constrained or full native inference.

\begin{table}[ht!]
\centering
\small
\caption{Historical shape-tied control on the first $N$ DeepSeek operations,
16 MiB global memory. Columns are LP variables; objectives use the original
cycle convention. Agreement on these cases is not a general equivalence proof.}
\label{tab:agg_validation}
\begin{tabular}{@{}rrrrccc@{}}
\toprule
$N$ & $S$ & cols (unagg.) & cols (agg.) & unaggregated & aggregated & error \\
\midrule
16  & 10 & 9{,}729   & 3{,}981 & $3.848169063\times10^{10}$ & $3.848169063\times10^{10}$ & $0.0000\%$ \\
32  & 10 & 37{,}377  & 3{,}981 & $1.027407872\times10^{11}$ & $1.027407872\times10^{11}$ & $0.0000\%$ \\
64  & 13 & 146{,}433 & 6{,}540 & $1.397050573\times10^{11}$ & $1.397050573\times10^{11}$ & $0.0000\%$ \\
128 & 13 & 579{,}585 & 6{,}540 & $2.088766669\times10^{11}$ & $2.088766669\times10^{11}$ & $0.0000\%$ \\
\bottomrule
\end{tabular}
\end{table}

\paragraph{Current scalability path.}
Section~\ref{sec:ocs} instead uses structured temporal operators and compact
working supports, retaining the source identity and checking omitted-domain
obligations separately. Shared hardware is global, but independent operation
choices remain independent. A working certificate applies to its identified
support; only a complete-domain check can extend a lower bound to the source.
